\documentclass[11pt]{article}

\usepackage[margin=1in]{geometry}
\usepackage{amsmath,amssymb,amsthm,mathtools}
\usepackage{enumitem}
\usepackage[hidelinks]{hyperref}
\usepackage{microtype}

\allowdisplaybreaks
\setlist[itemize]{topsep=2pt,itemsep=1pt,leftmargin=1.5em}
\setlist[enumerate]{topsep=2pt,itemsep=1pt,leftmargin=1.6em}

\newtheorem{theorem}{Theorem}[section]
\newtheorem{proposition}[theorem]{Proposition}
\newtheorem{lemma}[theorem]{Lemma}
\newtheorem{corollary}[theorem]{Corollary}
\theoremstyle{definition}
\newtheorem{definition}[theorem]{Definition}
\theoremstyle{remark}
\newtheorem{remark}[theorem]{Remark}

\newcommand{\nuu}{\nu_2}
\newcommand{\Ecal}{\mathcal E}
\newcommand{\Rcal}{\mathcal R}
\newcommand{\V}{\mathcal V}
\newcommand{\GD}{G_D}
\newcommand{\CD}{C_D}
\newcommand{\bdD}{\partial_D}
\newcommand{\depth}{\operatorname{depth}}
\newcommand{\dist}{\operatorname{dist}}
\newcommand{\oddrep}{\operatorname{oddrep}}
\newcommand{\Pre}{\mathsf{Pre}}
\newcommand{\Seed}{\operatorname{Seed}}
\newcommand{\Pair}{\mathsf P}
\newcommand{\Sfam}{\mathcal S}
\newcommand{\eps}{\varepsilon}

\title{Adaptive Valuation Graphs, Binomial Distance Layers, and Weighted Histories\\
\large A finite graph-theoretic structure from the \texorpdfstring{$3n+1$}{3n+1} valuation rule}
\author{}
\date{Fourth draft --- April 2026}

\begin{document}
\maketitle

\begin{abstract}
For each integer $D\ge 1$ we define a finite directed graph $G_D$ on odd residue classes modulo powers of two up to depth $D$.  The graph is adaptive: a class whose two-adic valuation data do not yet determine a positive-depth transition refines by one binary digit, while every other class carries the valuation-labeled transition determined by that data.  We prove that $G_D$ has a unique nontrivial strongly connected component $C_D$, give an explicit description of it, and show that $|C_D|=D(D-2)$ for $D\ge3$.  We also prove that valuation-edge preimages of a target of depth $e$ are naturally indexed by pairs $e<i<j\le D$.  This pair indexing yields a canonical history coordinate for vertices at positive distance from $C_D$.  After quotienting by the top-bit involution, the depth-$d$, distance-$k$ layer is naturally bijective with
\[
\Sfam_{d,k}=\{\Sigma\subseteq[d]: d\in\Sigma,\ |\Sigma|=2k+2\},
\]
and hence the oriented layer has cardinality
\[
\bigl|\{v:\depth(v)=d,\ R_D(v)=k\}\bigr|=2\binom{d-1}{2k+1}\qquad(k\ge1).
\]
Finally, when valuation edges are weighted by $t\log_2 3-(t+\rho)$, the total weight of a positive-distance chain factors through the same quotient history as an explicit alternating-gap functional.  The Collatz map motivates the local valuation rule and this particular weight, but the results are finite graph theory, weighted digraphs, and enumerative combinatorics.
\end{abstract}

\section{Introduction}

This paper studies a finite directed graph built from odd residue classes modulo powers of two.  The graph is not drawn at one fixed modulus.  Instead, it refines a residue class only when the currently visible two-adic valuation data are insufficient to define a positive-depth transition.  This adaptive rule produces a finite graph with a rigid structure: a single recurrent core, a small truncation boundary, pair-indexed valuation preimages, and binomial distance layers.  The same history coordinate also linearizes a natural additive weight on valuation chains.

The construction comes from the following two-adic valuation rule.  For an odd integer $X$, write
\[
X+1=2^tq,\qquad q\text{ odd},
\]
and then
\[
3^tq-1=2^\rho Y,\qquad Y\text{ odd}.
\]
The finite graph asks what can be determined from the residue class of $X$ modulo $2^d$.  If the class already determines the valuation data $(t,\rho)$ up to a positive target depth, the graph follows the induced transition.  If not, it asks for one more binary digit.

The main point is that the theorems below are not integer-orbit statements.  The graph and its combinatorics are the subject.  Collatz is the source of the valuation rule.  This distinction is important because the existing literature contains many graph, tree, $2$-adic, and probabilistic models of the $3x+1$ problem, but most of them are designed to encode integer or $2$-adic dynamics directly.  The present paper instead isolates a finite adaptive residue graph whose structural statements can be proved internally.

\section{Historical background and prior work}\label{sec:related}

The $3x+1$ problem has a large literature, and several strands are especially relevant to the present construction.  General surveys and bibliographies are indispensable for orientation: Lagarias' classical survey gives the standard formulation, early history, and many equivalent or generalized viewpoints \cite{Lagarias1985}; the later edited volume collects major expository and research articles on the problem \cite{Lagarias2010}; and the annotated bibliographies of Lagarias record the breadth of the literature through the first decades of systematic work \cite{LagariasAnnotated1999,LagariasAnnotated2009}.  Wirsching's monograph develops the $3n+1$ map as a dynamical system and treats predecessor sets, counting functions, and asymptotic questions in detail \cite{Wirsching1998}.  We use this literature mainly as context: the present results do not attempt to prove convergence of integer trajectories, but they do use a local valuation rule that comes from the odd-to-odd Syracuse dynamics.

A first line of prior work concerns stopping times and density theorems.  Terras \cite{Terras1976} and Everett \cite{Everett1977} proved that almost all positive integers eventually fall below their starting value.  This was sharpened by Allouche \cite{Allouche1979} and Korec \cite{Korec1994}.  More recently, Tao proved that almost all Collatz orbits attain almost bounded values, in logarithmic density, using an approximate-transport argument for Syracuse random variables and estimates involving a skew random walk on $3$-adic cyclic groups \cite{Tao2022}.  These results study integer orbits statistically.  Our weighted history functional is much more elementary: it assigns to a finite valuation edge the local logarithmic drift $t\log_2 3-(t+\rho)$ and proves that, on positive-distance layers of $G_D$, this drift factors through a binomial history coordinate.

A second line concerns inverse trees and predecessor sets.  The inverse $3x+1$ tree is already implicit in early work and was developed systematically in later counting and density results, including work of Applegate and Lagarias on the distribution of $3x+1$ trees and lower bounds for total stopping times \cite{ApplegateLagarias1995,ApplegateLagarias2003}.  Wirsching's book places predecessor sets at the center of the dynamical analysis \cite{Wirsching1998}.  Our preimage theorem is related in spirit but not in object: instead of enumerating integer predecessors of a fixed integer, it enumerates valuation-edge preimages of a residue vertex inside the finite adaptive graph.  The resulting index set is the elementary pair set $e<i<j\le D$, and iterating these pairs gives the subset histories in Theorem~\ref{thm:bijection}.

A third line is $2$-adic and symbolic.  Bernstein gave a non-iterative $2$-adic formulation of the conjecture \cite{Bernstein1994}.  Bernstein and Lagarias then studied the $3x+1$ conjugacy map $\Phi$ on the $2$-adic integers, conjugating the shift map to the $3x+1$ map and analyzing the induced permutations modulo $2^n$ \cite{BernsteinLagarias1996}.  Later work studied automorphisms and autoconjugacies of the $3x+1$ map, including the nontrivial autoconjugacy of Monks and Yazinski \cite{MonksYazinski2004}, and parity-vector viewpoints such as Rozier's treatment of parity sequences on $\mathbb Z_2$ \cite{Rozier2019}.  These works show that parity and $2$-adic coding are not auxiliary bookkeeping; they are central structure.  The present graph should be viewed in that same broad family, but with a different cutoff principle: it keeps vertices at varying depths and refines only when the valuation data have not yet determined a positive-depth transition.

A fourth line is explicitly graph-theoretic.  The ordinary Collatz graph is the directed graph on integers obtained by drawing arrows according to the Collatz rule; inverse versions produce rooted predecessor trees.  Laarhoven and de Weger studied congruence-class variants of the Collatz graph and proved that, for moduli that are powers of two, certain finite graphs are isomorphic to binary de Bruijn graphs; they also related the corresponding infinite graph isomorphism to the Bernstein--Lagarias conjugacy map \cite{LaarhovenDeWeger2013}.  This is the closest existing point of comparison.  Their graphs are fixed-modulus congruence graphs.  The graph $G_D$ here is adaptive: its vertices range over all depths $1\le d\le D$, and the edge rule either refines a class or follows a valuation-determined transition to a lower depth.  Consequently the main structural questions are not de Bruijn regularity at fixed modulus, but the recurrent core, the truncation boundary, pair-indexed valuation preimages, and binomial distance layers.

Finally, the combinatorial side of this paper is deliberately elementary.  We use standard directed-graph notions such as strongly connected components and distances in a digraph; see, for example, Diestel \cite{DiestelGraphTheory}.  The layer counts are ordinary binomial counts, and the quotient histories are subsets of $[d]$ of prescribed cardinality, in the spirit of basic enumerative combinatorics \cite{StanleyEC1}.  The novelty is not a new enumerative method, but the fact that the adaptive valuation graph produces these binomial layers canonically.

\subsection*{Position of this paper}

The contribution of the paper can be summarized as follows.  Starting from a local Syracuse valuation rule, we build a finite adaptive residue graph and prove a complete structural description of its positive-distance part.  The result sits between three familiar frameworks: residue dynamics modulo powers of two, $2$-adic symbolic codings, and inverse/predecessor graphs.  Unlike integer-orbit approaches, it makes no claim that Collatz trajectories converge.  Unlike fixed-modulus residue graphs, it allows the modulus depth to change under the edge rule.  Unlike predecessor-tree enumeration, it indexes finite graph preimages by pairs of depth coordinates.  This is why the main theorem is graph-theoretic and enumerative rather than a convergence theorem.

\section{The graph}\label{sec:graph}

For $d\ge1$, let $V_d$ be the set of odd residues modulo $2^d$, represented by odd integers $1\le r<2^d$.  A vertex of $G_D$ is a pair $(r,d)$ with $1\le d\le D$ and $r\in V_d$; write $\depth(r,d)=d$.  Thus
\[
V(G_D)=\{(r,d):1\le d\le D,\ r\equiv1\pmod2,\ 1\le r<2^d\}.
\]
If $x$ is odd modulo $2^m$, let $\oddrep_m(x)$ denote its canonical representative in $\{1,3,\ldots,2^m-1\}$.

\subsection{Exceptional residues and valuation edges}

For $1\le t<d$, define
\[
r_{d,t}\equiv 2^t(3^t)^{-1}-1\pmod {2^d},
\]
where the inverse is taken modulo $2^{d-t}$ and the odd representative modulo $2^d$ is chosen.  Also set
\[
r_{d,d}=2^d-1,
\qquad
\Ecal_d=\{r_{d,1},\ldots,r_{d,d}\}.
\]
The residues in $\Ecal_d$ are distinct, since $\nuu(r_{d,t}+1)=t$ for $t<d$ and $\nuu(r_{d,d}+1)=d$.  We call them exceptional.  All other residues at depth $d$ are regular.

The edges of $G_D$ are as follows.

If $(r,d)$ is exceptional and $d<D$, draw the two refinement edges
\[
(r,d)\to(r,d+1),\qquad (r,d)\to(r+2^d,d+1).
\]
If $d=D$, no refinement edge is drawn.

If $(r,d)$ is regular, put
\[
t=\nuu(r+1),\qquad q=\frac{r+1}{2^t},\qquad
\rho=\nuu(3^tq-1),\qquad e=d-t-\rho.
\]
Lemma~\ref{lem:local} below implies $e\ge1$.  Set
\[
y=\frac{3^tq-1}{2^\rho},\qquad s=\oddrep_e(y),
\]
and draw the valuation edge
\[
(r,d)\to(s,e).
\]
We denote this partial map by $F(r,d)=(s,e)$.

The truncation boundary is
\[
\bdD=\{(r,D):r\in\Ecal_D\}\cup\{(2^{D-1}-1,D-1)\}.
\]
The second term is the maximal exceptional vertex one level below the cutoff.

\begin{lemma}[Local valuation algebra]\label{lem:local}
The following hold.
\begin{enumerate}
\item A regular vertex has a valid valuation edge: $e=d-t-\rho\ge1$.
\item Every exceptional vertex at depth $d>1$ reduces to an exceptional vertex at depth $d-1$.
\item If $1\le t<d$, then the two children of $(r_{d,t},d)$ consist of one exceptional child of type $t$ and one regular child whose valuation edge has target $(1,1)$.
\item The two children of the maximal exceptional vertex $(2^d-1,d)$ are exceptional of types $d$ and $d+1$ at depth $d+1$.
\end{enumerate}
\end{lemma}

\begin{proof}
For a regular vertex, $e\le0$ would mean
\[
3^tq\equiv1\pmod {2^{d-t}},
\]
so $r=2^tq-1$ would be the exceptional residue of type $t$, a contradiction.

The parent statement follows immediately from the congruence defining $r_{d,t}$; the case $t=d$ reduces to the maximal exceptional residue at depth $d-1$.

For the child statement, write $r=r_{d,t}$ with $t<d$ and $q=(r+1)/2^t$.  The two depth-$(d+1)$ lifts replace $q$ by $q$ and $q+2^{d-t}$.  Since
\[
3^t(q+2^{d-t})-1=(3^tq-1)+3^t2^{d-t},
\]
exactly one lift satisfies the stronger congruence modulo $2^{d+1-t}$; this is the exceptional child of type $t$.  The other has exact valuation $\rho=d-t$, hence target depth $(d+1)-t-(d-t)=1$, so it maps to $(1,1)$.  Finally, the children of $2^d-1$ are $2^d-1$ and $2^{d+1}-1$, which are the exceptional residues of types $d$ and $d+1$ at depth $d+1$.
\end{proof}

A regular vertex is a \emph{direct return} if its valuation edge lands at $(1,1)$.  Let $\Rcal_d$ be the set of direct returns at depth $d$.  Lemma~\ref{lem:local} gives
\[
|\Rcal_d|=\max(d-2,0),
\]
and for $d\ge3$ these are exactly the regular children of the non-maximal exceptional vertices at depth $d-1$.

\section{The recurrent core}\label{sec:core}

For $D\ge3$, define
\[
K_D=\left(\{(r,d):r\in\Ecal_d,\ 1\le d\le D\}\cup\bigcup_{d=1}^D\Rcal_d\right)\setminus\bdD.
\]
Thus $K_D$ consists of the exceptional vertices and direct returns, except for the truncation boundary.

\begin{theorem}[Core theorem]\label{thm:core}
For $D\ge3$, $G_D$ has a unique nontrivial strongly connected component $C_D$, and
\[
C_D=K_D.
\]
Moreover
\[
|C_D|=D(D-2),
\]
and a vertex reaches $C_D$ if and only if it is not in $\bdD$.
\end{theorem}

\begin{proof}
First, $K_D$ is strongly connected.  Every direct return maps to $(1,1)$.  If an exceptional vertex in $K_D$ has non-maximal type, Lemma~\ref{lem:local} gives a direct-return child and hence a path to $(1,1)$.  If it is maximal, it is not in the boundary, so its depth is at most $D-2$; one refinement step reaches a non-maximal exceptional child, and then the previous case applies.  Conversely, every exceptional vertex is reachable from $(1,1)$ by repeatedly taking exceptional parents, and every direct return is a refinement child of an exceptional vertex.  Hence $(1,1)$ reaches all of $K_D$.

Now let $v\notin K_D\cup\bdD$.  Then $v$ is regular and not a direct return.  Its unique valuation edge strictly lowers depth, since $t\ge1$ and $\rho\ge1$.  Iterating therefore reaches either $K_D$ or the boundary.  It cannot first hit the boundary: every valuation edge from a source of depth at most $D$ has target depth at most $D-2$, while the boundary lies only in depths $D-1$ and $D$.  Hence every non-boundary vertex reaches $K_D$.

The boundary does not reach $K_D$: terminal exceptional vertices have no outgoing refinement edges, and the extra boundary vertex $(2^{D-1}-1,D-1)$ refines only to terminal exceptional vertices.  Finally, outside $K_D$ any path either enters the boundary trap or follows valuation edges of strictly decreasing depth, so no further directed cycle is possible.  Thus $K_D$ is the unique nontrivial strongly connected component.

For the count,
\[
\sum_{d=1}^D |\Ecal_d|=\sum_{d=1}^D d=\frac{D(D+1)}2,
\]
and
\[
\sum_{d=1}^D |\Rcal_d|=\sum_{d=3}^D(d-2)=\frac{(D-1)(D-2)}2.
\]
Their union has size $D^2-D+1$.  Removing the $D$ terminal exceptional vertices and the additional boundary vertex gives
\[
|C_D|=D^2-D+1-(D+1)=D(D-2).
\]
\end{proof}

Let
\[
R_D(v)=\dist(v,C_D)
\]
when $v$ reaches $C_D$, and $R_D(v)=\infty$ otherwise.  Theorem~\ref{thm:core} says that $R_D(v)<\infty$ exactly off the boundary.  If $R_D(v)=k>0$, then $v$ is regular and not a direct return, and repeated valuation edges give a unique chain
\[
v=v_k\to v_{k-1}\to\cdots\to v_0\in C_D.
\]

\section{Preimages and pair indexing}\label{sec:preimages}

The next theorem is the basic enumerative mechanism.  It is target-first: fix a target and solve for all valuation-edge sources.

\begin{theorem}[Preimage-pair theorem]\label{thm:preimage}
Fix a target vertex $u=(s,e)\in V(G_D)$.  The valuation-edge preimages of $u$ are naturally indexed by pairs
\[
e<i<j\le D.
\]
For the pair $(i,j)$, set
\[
t=i-e,
\qquad
\rho=j-i,
\qquad
 d=j.
\]
Then the unique source is obtained by solving
\[
q\equiv3^{-t}(1+2^\rho s)\pmod {2^{e+\rho}}
\]
and setting
\[
\Pre_{i,j}(s,e)=\bigl(\oddrep_j(2^tq-1),j\bigr).
\]
Consequently $|F^{-1}(s,e)|=\binom{D-e}{2}$.
\end{theorem}

\begin{proof}
If a valuation edge has source depth $d$ and target depth $e$, then
\[
d=e+t+\rho.
\]
Thus $i=e+t$ and $j=d$ satisfy $e<i<j\le D$.

Conversely, fix $e<i<j\le D$ and put $t=i-e$, $\rho=j-i$.  A source with target residue $s$ must satisfy
\[
3^tq-1\equiv2^\rho s\pmod {2^{e+\rho}},
\]
or equivalently
\[
q\equiv3^{-t}(1+2^\rho s)\pmod {2^{e+\rho}}.
\]
This congruence has a unique odd solution modulo $2^{e+\rho}$.  With $r=\oddrep_j(2^tq-1)$, one has $\nuu(r+1)=t$.  Since $s$ is odd, the valuation of $3^tq-1$ is exactly $\rho$, and division by $2^\rho$ gives target residue $s$ modulo $2^e$.  The source is regular because $e\ge1$.  Uniqueness follows from the same congruence for $q$.
\end{proof}

The pair attached to a valuation edge $v=(r,d)\to(s,e)$ is therefore
\[
\Pair(v)=\{e+\nuu(r+1),d\}.
\]
Each preimage step adds two depths: an intermediate valuation depth and the source depth.

\section{Histories and binomial layers}\label{sec:histories}

For $d\ge2$, let $\tau_d$ be the top-bit involution
\[
\tau_d(r,d)=
\begin{cases}
(r+2^{d-1},d),&r<2^{d-1},\\
(r-2^{d-1},d),&r>2^{d-1}.
\end{cases}
\]
It exchanges the two depth-$d$ lifts of the same residue modulo $2^{d-1}$.

\begin{proposition}[Stable core seeds]\label{prop:seeds}
For $2\le d\le D-2$, the quotient $(C_D\cap V_d)/\tau_d$ is canonically labeled by the two-element sets
\[
\{a,d\},\qquad 1\le a<d.
\]
There are two oriented core vertices above each label.
\end{proposition}

\begin{proof}
For $1\le a\le d-2$, the two children of the exceptional vertex of type $a$ at depth $d-1$ consist of the exceptional vertex of type $a$ at depth $d$ and its direct-return sibling.  These two children differ by $2^{d-1}$, so they form one $\tau_d$-orbit.  The remaining orbit comes from the two children of the maximal exceptional vertex at depth $d-1$, namely the exceptional vertices of types $d-1$ and $d$.  Since $d\le D-2$, no boundary vertex is involved.  These $d-1$ two-element orbits exhaust $C_D\cap V_d$.
\end{proof}

If $v=v_k\to\cdots\to v_0\in C_D$ is the unique chain for a vertex of positive finite distance, write $v_m=(r_m,d_m)$.  The endpoint $v_0$ has stable depth: $2\le d_0\le D-2$, because the final edge is not a direct return and every valuation edge lowers depth by at least two.  Let
\[
\Seed(v_0)=\{a,d_0\}
\]
be its core seed label from Proposition~\ref{prop:seeds}.  Define the projected history
\[
\Sigma(v)=\Seed(v_0)\cup\bigcup_{m=1}^k\{d_{m-1}+\nuu(r_m+1),d_m\}.
\]
For a stable core vertex itself, set $\Sigma(v_0)=\Seed(v_0)$.

\begin{lemma}[History intervals]\label{lem:historyintervals}
If $v$ has finite distance $k$ and depth $d$, then the union defining $\Sigma(v)$ is disjoint.  For $k\ge1$,
\[
\Sigma(v)\subseteq[d],\qquad d\in\Sigma(v),\qquad |\Sigma(v)|=2k+2.
\]
\end{lemma}

\begin{proof}
For the edge $v_m\to v_{m-1}$, put
\[
t_m=\nuu(r_m+1),\qquad \rho_m=\nuu(3^{t_m}q_m-1).
\]
Then
\[
d_m=d_{m-1}+t_m+\rho_m,
\]
so
\[
\{d_{m-1}+t_m,d_m\}\subset(d_{m-1},d_m].
\]
The seed lies in $[1,d_0]$.  Inducting along the chain, each preimage step appends two new elements in the next interval $(d_{m-1},d_m]$.  Hence the union is disjoint, lies in $[d]$, contains the final depth $d$, and has cardinality $2+2k$.
\end{proof}

For $k\ge1$, define
\[
\Sfam_{d,k}=\{\Sigma\subseteq[d]:d\in\Sigma,\ |\Sigma|=2k+2\}.
\]

\begin{theorem}[History bijection]\label{thm:bijection}
For $D\ge3$, $k\ge1$, and $1\le d\le D$, the map
\[
v\longmapsto (\eps(v),\Sigma(v))
\]
is a natural bijection
\[
\{v:\depth(v)=d,\ R_D(v)=k\}
\longrightarrow
\{+,-\}\times\Sfam_{d,k},
\]
where $\eps(v)$ is the orientation of the terminal core seed.  Consequently
\[
\bigl|\{v:\depth(v)=d,\ R_D(v)=k\}\bigr|=2\binom{d-1}{2k+1}.
\]
\end{theorem}

\begin{proof}
Lemma~\ref{lem:historyintervals} shows that the map lands in $\{+,-\}\times\Sfam_{d,k}$.

To reconstruct, take $(\eps,\Sigma)$ with
\[
\Sigma=\{\sigma_1<\sigma_2<\cdots<\sigma_{2k+2}=d\}.
\]
Start with the oriented core vertex labeled by $\{\sigma_1,\sigma_2\}$ and orientation $\eps$.  At step $m=1,\ldots,k$, use Theorem~\ref{thm:preimage} to take the unique valuation-edge preimage of the current target of depth $\sigma_{2m}$ indexed by the pair
\[
\sigma_{2m}<\sigma_{2m+1}<\sigma_{2m+2}.
\]
This constructs a chain whose final source has depth $d$.

Each reconstructed edge has target depth at least $2$, so no reconstructed source is a direct return; hence the first time the chain reaches the core is its prescribed endpoint.  Thus the final source has distance $k$.  The construction is inverse to the history map because Theorem~\ref{thm:preimage} gives a unique source for each target and pair.

The count follows from
\[
|\Sfam_{d,k}|=\binom{d-1}{2k+1},
\]
since $d$ is forced to lie in $\Sigma$.
\end{proof}

The top-bit quotient removes the two orientations.  The needed equivariance is elementary: if a valuation edge has target depth $e\ge2$, then replacing the source by its top-bit flip changes $q$ to $q\pm2^{d-1-t}$, preserves the exact valuations $(t,\rho)$, and flips the top bit of the target.  Thus valuation chains commute with $\tau$ as long as no direct return is involved.  Positive-distance chains have this property.

\begin{corollary}[Quotient layers]\label{cor:quotient}
For $k\ge1$ and $d\ge2$,
\[
\{v:\depth(v)=d,\ R_D(v)=k\}/\tau_d
\cong
\Sfam_{d,k}.
\]
Equivalently, summing over $d\le D$,
\[
\bigl|\{v:R_D(v)=k\}\bigr|=2\binom{D}{2k+2}.
\]
\end{corollary}


\section{Weighted histories}\label{sec:weighted}

The history coordinate is not only a counting device.  It also records a natural additive statistic on valuation chains.  In this section we regard $G_D$ as a weighted directed graph by assigning a real weight to each valuation edge from its valuation label.

Put
\[
\alpha=\log_2 3.
\]
If $v=(r,d)$ is regular and $F(v)=u=(s,e)$, write
\[
t=\nuu(r+1),\qquad q=\frac{r+1}{2^t},\qquad \rho=\nuu(3^tq-1).
\]
Thus $d=e+t+\rho$.  Define the edge weight
\[
\omega(v,u)=t\alpha-(t+\rho)=t\alpha-(d-e).
\]
Equivalently, if the pair of the edge is
\[
\Pair(v)=\{i,j\},\qquad e<i<j=d,
\]
then
\[
\omega(v,u)=(i-e)\alpha-(j-e).
\]

Now let $v$ have positive finite distance $R_D(v)=k$, and write its valuation chain as
\[
v=v_k\to v_{k-1}\to\cdots\to v_0\in C_D.
\]
The total chain weight is
\[
\Omega(v)=\sum_{m=1}^k \omega(v_m,v_{m-1}).
\]
The endpoint core vertex carries no weight in this definition; only the valuation edges in the distance chain are counted.

For a history
\[
\Sigma=\{\sigma_1<\sigma_2<\cdots<\sigma_{2k+2}=d\}\in\Sfam_{d,k},
\]
define the alternating-gap sum
\[
T(\Sigma)=\sum_{m=1}^k(\sigma_{2m+1}-\sigma_{2m})
\]
and the history weight
\[
\Omega(\Sigma)=T(\Sigma)\alpha-(d-\sigma_2).
\]

\begin{theorem}[Weight factorization through histories]\label{thm:weightfactor}
Let $D\ge3$, $k\ge1$, and let $v$ be a vertex with $R_D(v)=k$ and $\depth(v)=d$.  Then
\[
\Omega(v)=\Omega(\Sigma(v)).
\]
Consequently the total chain weight depends only on the quotient history $\Sigma(v)$ and is independent of the orientation $\eps(v)$.
\end{theorem}

\begin{proof}
Write the distance chain as
\[
v=v_k\to v_{k-1}\to\cdots\to v_0\in C_D,
\]
and put $v_m=(r_m,d_m)$.  For the edge $v_m\to v_{m-1}$, let
\[
t_m=\nuu(r_m+1),\qquad q_m=\frac{r_m+1}{2^{t_m}},\qquad \rho_m=\nuu(3^{t_m}q_m-1).
\]
By the definition of the history coordinate, if
\[
\Sigma(v)=\{\sigma_1<\sigma_2<\cdots<\sigma_{2k+2}=d\},
\]
then the $m$th valuation preimage step has
\[
d_{m-1}=\sigma_{2m},\qquad d_{m-1}+t_m=\sigma_{2m+1},\qquad d_m=\sigma_{2m+2}.
\]
Therefore
\[
t_m=\sigma_{2m+1}-\sigma_{2m},
\qquad
 t_m+\rho_m=d_m-d_{m-1}=\sigma_{2m+2}-\sigma_{2m}.
\]
Summing the edge weights gives
\begin{align*}
\Omega(v)
&=\sum_{m=1}^k\bigl(t_m\alpha-(t_m+\rho_m)\bigr)\\
&=\alpha\sum_{m=1}^k(\sigma_{2m+1}-\sigma_{2m})
  -\sum_{m=1}^k(\sigma_{2m+2}-\sigma_{2m})\\
&=T(\Sigma(v))\alpha-(\sigma_{2k+2}-\sigma_2)\\
&=T(\Sigma(v))\alpha-(d-\sigma_2).
\end{align*}
This is exactly $\Omega(\Sigma(v))$.  Since $\Sigma(v)$ is unchanged by switching the top-bit orientation, the total chain weight is independent of $\eps(v)$.
\end{proof}

\begin{corollary}[Weighted quotient layers]\label{cor:weightedquotient}
For $k\ge1$ and $d\ge2$, the quotient bijection
\[
\{v:\depth(v)=d,\ R_D(v)=k\}/\tau_d\cong\Sfam_{d,k}
\]
is weight-preserving when the left side is weighted by total chain weight and $\Sfam_{d,k}$ is weighted by $\Omega(\Sigma)$.
\end{corollary}

\begin{remark}
The statistic $T(\Sigma)$ selects the odd-numbered gaps after the core seed depth $\sigma_2$, while $d-\sigma_2$ is the total depth gained after the seed.  Thus the weight is an alternating-gap functional on the same subsets that index the binomial distance layers.  No information from the orientation coordinate is needed.
\end{remark}

\section{Relation to Collatz}\label{sec:collatz}

The rule used to define $G_D$ is extracted from a residue-class model of the odd Collatz/Syracuse dynamics.  In that model one studies the odd-to-odd valuation update
\[
X+1=2^tq,
\qquad
3^tq-1=2^\rho Y.
\]
This origin explains why the formulas occur.  It also explains the particular edge weight $t\log_2 3-(t+\rho)$ used in Section~\ref{sec:weighted}.  Neither point is needed for the finite graph theorems: the graph, its layers, and its weighted histories are the objects of study here.

The statements proved here do not assert that integer Collatz orbits reach $1$.  A vertex $(r,d)$ is a residue class, not an integer.  The integer $1$ is represented by an infinite compatible tower of residue classes, and integer dynamics require lift data beyond any fixed $G_D$.  What the graph supplies is a finite, classifiable combinatorial object: adaptive two-adic refinement produces a recurrent core, pair-indexed valuation preimages, a top-bit symmetry, binomial distance layers, and a history-level weight functional.


\begin{thebibliography}{99}

\bibitem{Allouche1979}
J.-P. Allouche.
Sur la conjecture de ``Syracuse--Kakutani--Collatz''.
\emph{Seminaire de Theorie des Nombres de Bordeaux}, exp. no.~9, 1978--1979.

\bibitem{ApplegateLagarias1995}
D. Applegate and J. C. Lagarias.
The distribution of $3x+1$ trees.
\emph{Experimental Mathematics} 4 (1995), no.~3, 193--209.

\bibitem{ApplegateLagarias2003}
D. Applegate and J. C. Lagarias.
Lower bounds for the total stopping time of $3x+1$ iterates.
\emph{Mathematics of Computation} 72 (2003), no.~242, 1035--1049.

\bibitem{Bernstein1994}
D. J. Bernstein.
A non-iterative $2$-adic statement of the $3n+1$ conjecture.
\emph{Proceedings of the American Mathematical Society} 121 (1994), no.~2, 405--408.

\bibitem{BernsteinLagarias1996}
D. J. Bernstein and J. C. Lagarias.
The $3x+1$ conjugacy map.
\emph{Canadian Journal of Mathematics} 48 (1996), no.~6, 1154--1169.

\bibitem{Conway1972}
J. H. Conway.
Unpredictable iterations.
In \emph{Proceedings of the 1972 Number Theory Conference}, University of Colorado, Boulder, 1972, 49--52.

\bibitem{DiestelGraphTheory}
R. Diestel.
\emph{Graph Theory}.
Graduate Texts in Mathematics 173, 5th ed., Springer, 2017.

\bibitem{Everett1977}
C. J. Everett.
Iteration of the number-theoretic function $f(2n)=n$, $f(2n+1)=3n+2$.
\emph{Advances in Mathematics} 25 (1977), no.~1, 42--45.

\bibitem{Garner1981}
L. E. Garner.
On the Collatz $3n+1$ algorithm.
\emph{Proceedings of the American Mathematical Society} 82 (1981), no.~1, 19--22.

\bibitem{Gouvea1997}
F. Q. Gouvea.
\emph{$p$-adic Numbers: An Introduction}.
2nd ed., Universitext, Springer, 1997.

\bibitem{Korec1994}
I. Korec.
A density estimate for the $3x+1$ problem.
\emph{Mathematica Slovaca} 44 (1994), no.~1, 85--89.

\bibitem{Lagarias1985}
J. C. Lagarias.
The $3x+1$ problem and its generalizations.
\emph{The American Mathematical Monthly} 92 (1985), no.~1, 3--23.

\bibitem{LagariasAnnotated1999}
J. C. Lagarias.
The $3x+1$ problem: an annotated bibliography (1963--1999).
Preprint, arXiv:math/0309224.

\bibitem{LagariasAnnotated2009}
J. C. Lagarias.
The $3x+1$ problem: an annotated bibliography, II (2000--2009).
Preprint, arXiv:math/0608208.

\bibitem{Lagarias2010}
J. C. Lagarias, editor.
\emph{The Ultimate Challenge: The $3x+1$ Problem}.
American Mathematical Society, Providence, RI, 2010.

\bibitem{LaarhovenDeWeger2013}
T. Laarhoven and B. de Weger.
The Collatz conjecture and de Bruijn graphs.
\emph{Indagationes Mathematicae} 24 (2013), no.~4, 971--983.

\bibitem{MonksYazinski2004}
K. G. Monks and J. Yazinski.
The autoconjugacy of the $3x+1$ function.
\emph{Discrete Mathematics} 275 (2004), no.~1--3, 219--236.

\bibitem{Rozier2019}
O. Rozier.
Parity sequences of the $3x+1$ map on the $2$-adic integers and Euclidean embedding.
\emph{Integers} 19 (2019), Paper A8.

\bibitem{StanleyEC1}
R. P. Stanley.
\emph{Enumerative Combinatorics, Volume 1}.
2nd ed., Cambridge Studies in Advanced Mathematics 49, Cambridge University Press, 2011.

\bibitem{Tao2022}
T. Tao.
Almost all orbits of the Collatz map attain almost bounded values.
\emph{Forum of Mathematics, Pi} 10 (2022), e12.

\bibitem{Terras1976}
R. Terras.
A stopping time problem on the positive integers.
\emph{Acta Arithmetica} 30 (1976), 241--252.

\bibitem{Wirsching1998}
G. J. Wirsching.
\emph{The Dynamical System Generated by the $3n+1$ Function}.
Lecture Notes in Mathematics 1681, Springer, 1998.

\end{thebibliography}

\end{document}
